\section{The geometric witnesses at the two distinguished primes}

% BEGIN T79W1
\paragraph{T79W1. Thresholds for the original type-prime condition.}
Use the setup of F78D2 and its chosen parameter. Write
$b=L.\mathrm{sc}.b$, $\eta=L.\mathrm{etaSz}$ and $t=D.R.t$.
Let $C_T$ be the prescribed natural number with $C_T\le b$.
Take a twisted localization $i$ with $\ell_i=g_k$, $k\in\{0,1\}$,
and with the source's nonemptiness predicate $\mathrm{HG}(i)$.
Use the full point family $R_x=X.r_i(x)$ and its actual complement
$\rho_x=X.\rho_i(x)$. The point-family fields of F78D2 give
$\mathrm{GoodR}(R)$ and, at each $x\in Z(i)$, an identified
complement and semisimplicity of $\rho_x\oplus\bigoplus_{c=1}^t\chi_c$.
These assertions hold on $Z(i)$, including the subset $Z_{\rm good}(i)$.

Set $s_* =\lfloor3b\rfloor+2$ and $S_*=0$. The canonical hub has
$X.\eta=\lfloor\eta\rfloor$; the parameter inequalities give
$\eta+t\le b$, $b\ge1$, hence $X.\eta+t+1<s_*$. Also
$X.\eta\le2\eta$. For the other strict inequality put $k_0=D.R.k$ and $y=k_0 b$. Then $k_0\ge5$,
$C_T\le b\le y$, $s_*\le3b+2\le y$ and
$n_V+1\le2k_0+t\le3y$. The last rank bound follows from
$n_V=k_0-1+D.R.m+t$, $D.R.m+3\le k_0$ and $t\le b$. Consequently
\[
 C_T\big((s_*+S_*)(n_V+1)^2+(n_V+1)^3\big)
 \le36y^4<y^8\le g_k .
\]
Here $y\ge5$ and the last inequality is the type-prime size condition
of the very chosen parameter. These numbers work for both indices $k$.
% END T79W1

% BEGIN T79W2
\paragraph{T79W2. Conjugate signs and jump sets for every point.}
The prime $g_k$ splits completely in the Galois host $E$. Indeed the
source gives exactly $[E:\mathbb Q]$ primes above it; the sum of their
residue degrees times ramification indices is $[E:\mathbb Q]$,
so each local degree is one. Each $w\mid\ell_i$ therefore determines
a unique embedding $E\hookrightarrow\overline{\mathbb Q}_{\ell_i}$
inducing $w$, through $E_w=\mathbb Q_{\ell_i}$. Equivalently, two such
embeddings differ by the trivial decomposition group at $w$.
There is likewise just one embedding of the local field into the
prescribed coefficient extension. The embeddings for conjugate
places differ by precomposition with the host's complex conjugation.

Use the sign convention $e$ supplied once by $\mathrm{AdmDict}$.
Transport each embedding to $\mathbb C$ by the given coefficient
isomorphism, label it by $g_w$ in the source's enumeration, and
choose its associated archimedean place. Set $s_g(w)$ to the truth
value of $(g_w\in\Psi)\Leftrightarrow(e=\text{true})$.
Since $\Psi$ contains exactly one element of every CM-conjugate pair,
the two signs at a conjugate pair are opposite. Give primes outside
$\ell_i$ sign false. This proves the actual $\mathrm{SignOK}$.

The type-prime clause of $\mathrm{AdmJ}$ supplies de Rhamness and the
ambient archimedean multiset at $w$ for \emph{every} $R_x$.
Thus the dictionary and reading calculation of S79P3 applies to
the ambient $R_x\chi_1^{-1}$ here, before subtracting the character
block. Write $d_0=(t-1)/2$, $a_*=e_{j-1}$, and use the entries
$v_a,b_a$ at the chosen archimedean place. Take the integral members
of the lists
\[
 d_0+a_*+v_a,\qquad d_0+a_*-b_a,\qquad 0,\ldots,t-1,
\]
after reflecting each $z$ to $2d_0-z$ if the sign is false.
Call their set $U(w)$; outside $w\mid\ell_i$ put $U(w)=\varnothing$.
The singleton rank-one weight and the shift of the ambient
multiplicities in $\mathrm{AdmRead}$ give exactly these possible jumps,
with either sign, at all $x$ and all local embeddings.
The cardinality of $U(w)$ is at most $(D.R.k-1)+D.R.m+t=n_V$.

The middle list has $|a_*-b_a|\le\eta$ and
$|a_*-b_a|\ge(t+3)/2=d_0+2$, by the ordered archimedean gaps in
S79P4. It lies within $\lfloor\eta\rfloor+1$ of $d_0$ and
contains neither $-1$ nor $t$. The first list lies strictly above
$t+1$ before reflection, since $a_*>\eta+t+L.\mathrm{Sc}$,
$\eta\ge t$ and $t\ge1$. Its entries are at most $2b\le s_*$,
by the same size and width inequalities used in S79P4. Thus their
signed distances from $d_0$ lie in $(0,s_*+S_*-d_0]$, with
the positive direction set by $s_g(w)$. Reflection keeps both
exclusions and preserves the third list. The third list itself
is the allowed interval $[0,t)$. This proves every clause of
$\mathrm{HostJ}(i,R,X.\eta,s_*,S_*,s_g,U)$, including the
exclusions, independently of any point in $Z(i)$.
% END T79W2

% BEGIN T79W3
\paragraph{T79W3. The local parameter at the other type prime.}
At a place $v\in S_i$ over a distinguished prime and away from
$\ell_i$, and for every point $x$, there is a point representative
$P$ with local Weil--Deligne parameter of $R_x$ equal to that of
$P_v$ twisted by $(1-n_V)/2$. Here and below the equality is in
the original parameter space, after the specified coefficient map.
We recall how the base-change and attachment fields of the setup
supply this fact, since $\mathrm{GoodR}$ by itself prescribes
Frobenius polynomials only outside the bad set.

Choose $P$ above $x$ by point surjectivity, and its isobaric base
change by $\mathrm{adm1}$. Write its parts $(\pi,a,d)$, where $a$
is the cuspidal rank and $d$ the number of shifts; the base-change
clause supplies conjugate self-duality of each $\pi$. For each
$0\le j<d$ the required twist of the cuspidal parameter is $e/2$,
where $e=d-n_V-2j$. Indeed the Speh shifts $(d-1)/2-j$ must be
added to the common normalization $(1-n_V)/2$. At every complex
embedding the multiset of shifted cuspidal weights is a submultiset
of the regular algebraic infinitesimal character of $P$. Its
entries are distinct and belong to $\mathbb Z+(n_V-1)/2$.
Thus for each cuspidal weight $z$ we have $z+e/2\in\mathbb Z$.
The attachment of this half-integral twist uses exactly the supplied
cuspidal and rank-one fields of $\mathrm{adm2}$. If $e+a-1=2h$,
$\pi$ itself is regular algebraic in the rank-$a$ normalization.
Attach its representation and tensor by the character attached to
the trivial automorphic character with shift $h$. If
$e+a-1=2h+1$, twist $\pi$ first by the supplied conjugate-self-dual
character $\eta_T$ of archimedean weights $\pm1/2$.
The distinct weights now have the rank-$a$ algebraic parity.
Tensor its attachment by the rank-one attachment to
$\eta_T^{-1}$ with shift $h+1/2$; its weights with this shift are
integers. In both cases the characters cancel and the local shifts add
to $(1-a)/2+(e+a-1)/2=e/2$, by the given ties for local twists.
The constituents are semisimple and tensoring by a character preserves
semisimplicity. Their direct sum is continuous semisimple of rank
$\sum ad=n_V$, and its parameters at places away from $\ell_i$ are
the normalized direct sums specified by the base change.

At split primes outside the finite excluded set the Hecke-polynomial
tie at $x$ and the base-change tie at $P$ identify the Frobenius
characteristic polynomials of $R_x$ and of this sum. The split
Frobenius density and semisimple character fields of $\mathrm{adm5}$
now identify the two Galois representations: each characteristic
coefficient is continuous, so coincidence on the dense set extends
to every group element, and equality of all these polynomials gives
conjugacy. Here $R_x$ is extended to the algebraic closure of its
coefficient field; semisimplicity persists because the characteristic
is zero (the finite-dimensional semisimple image algebra is separable).
The place $v$ under consideration is split, so the isobaric
base-change parameter there is that of $P_v$. Transporting
the Weil--Deligne parameter under the conjugacy proves the assertion.
% END T79W3

% BEGIN T79W4
\paragraph{T79W4. Wild inertia and the specified tame generators.}
Every $v\in S_i\cup Q_i$ has residue characteristic $p_v\ne\ell_i$.
The native character $\chi_1$ is trivial on wild inertia there.
To see this, its $\psi_L$ factor has open kernel and $\ell_i$-power
order on each element. The image of the pro-$p_v$ wild group in its
discrete quotient is a finite $p_v$-group, hence is trivial.
For the remaining twist character at the auxiliary places its inertia
action is a power, possibly an inverse, of the prescribed character of
order $15$. The auxiliary norm is $1$ modulo $15$ so $p_v\nmid15$.
At places above $g_a$ the inertia order divides $g_a-1$, again
prime to $p_v$. Elsewhere away from $\ell_i$ that character is
unramified. In the finite-order cases the continuous image in the
finite group of roots of the indicated order is also a $p_v$-group,
and therefore trivial on wild inertia.

If $v\in Q_i$, or $v$ is the designated $\mathrm{pde}$ place
(or the source's $\mathrm{pde}'$ alternative if
$\ell_i=\mathrm{pde}$), $\mathrm{ConclusionG5}$ conjugates the
inertia matrices of each $R_x$ over the fraction field to a diagonal
matrix whose entries are $1$ and $\chi_1$. Thus they act as the
identity on wild inertia, also over the integral coefficient ring
by injectivity. By the definition of $S_i$, every remaining place
lies above a distinguished prime. The depth-zero field of the setup
makes the inertia polynomial of $P_v$ on every wild element
$(Y-1)^{n_V}$. T79W3 and the inertia polynomial tie show the same
for $R_x$: the normalization twist is unramified and the coefficient
inclusion is injective. For clarity, if $B$ is such a matrix then
$(B-I)^{n_V}=0$ by Cayley--Hamilton. Modulo any power of the
coefficient maximal ideal it has $\ell_i$-power order: expand
$(I+(B-I))^{\ell_i^b}$ and, for each $1\le j<n_V$, use
$v_{\ell_i}\binom{\ell_i^b}{j}\ge b-v_{\ell_i}(j)$.
Taking $b$ large makes every surviving coefficient zero in the
quotient. The wild image in that discrete quotient is a finite
$p_v$-group by continuity and compactness. Consequently the
matrix is the identity modulo every ideal power, hence the identity
in the separated coefficient DVR. This proves wild triviality for
all $x$.

The tame-generator field of the datum singles out $\gamma_0(i,v)$
and says that it and wild inertia generate a dense subgroup of
inertia. At $q\in Q_i$ the first local character is ramified
on inertia. If $\zeta(i,q)=\chi_1(\gamma_0(i,q))$ were one,
its closed kernel would contain that dense subgroup, by wild
triviality just proved, contrary to ramification. This gives the
generator, wild-action and nonidentity clauses of $\mathrm{SetSQ}$
for exactly $R=X.r_i$. Combining T79W1--T79W4 gives the type-prime
witness extracted as the twelfth component of
\texttt{split\_bound\_of\_params}, on the original hub and for
either prime at every localization covered by its quantifiers.
% END T79W4
